\chapter{Hypothesis Register}\label{app:hypotheses}
\Cref{tab:hyp} lists every confirmatory hypothesis with its margin, its recomputed result at both targets and its verdict. Three further hypotheses of the first pre-registration were not tested or not confirmatory: H2 (D against the controls B and C) and H5 (D against E) were gated on H1 and were therefore \emph{not tested}; H3 (gradient alignment) is a mechanism check that is reported in \cref{sec:h3}; H4 (stochastic gradient descent with momentum in a deeper network) lies outside the experiments of this report.

\begin{table}[h]
\centering\small
\caption{Confirmatory hypotheses. Intervals: $97.5\%$ for H1, H6, H7a, H7b and $98.3\%$ for H8--H10, paired bootstrap. A verdict needs both targets to agree (\cref{sec:stats}). The ``margin'' is the saving that must be exceeded for \emph{supported}.}
\label{tab:hyp}
\resizebox{\textwidth}{!}{%
\begin{tabular}{@{}l l l l l l@{}}
\toprule
 & Comparison & Margin & change to 0.97 (\%) & change to 0.975 (\%) & Verdict\\
\midrule
H1 & D vs.\ A & 15\% & $+36.9$ [$+26.4$, $+47.4$] & $+35.5$ [$+26.9$, $+44.4$] & refuted \\
H6 & B vs.\ A & 15\% & $+2.3$ [$-4.8$, $+10.4$] & $+1.7$ [$-5.0$, $+9.2$] & refuted \\
\midrule
H7a & S30 vs.\ A1 & 15\% & $-35.1$ [$-38.6$, $-31.5$] & $-38.4$ [$-41.2$, $-35.6$] & supported \\
H7b & S30 vs.\ B48 & 0\% & $-30.2$ [$-35.7$, $-24.3$] & $-35.5$ [$-39.8$, $-30.9$] & supported \\
\midrule
H8 & S20w0 vs.\ A1 & 15\% & $-48.5$ [$-52.1$, $-44.2$] & $-48.8$ [$-53.4$, $-43.8$] & supported \\
H9 & S20w0 vs.\ S30 & 10\% & $-22.2$ [$-29.1$, $-14.7$] & $-19.1$ [$-26.1$, $-11.7$] & supported \\
H10 & S20w0 vs.\ SBT & 15\% & $-46.5$ [$-52.3$, $-40.3$] & $-53.6$ [$-58.3$, $-48.2$] & supported
\\
\bottomrule
\end{tabular}}
\end{table}

\paragraph{Reading the table.} A negative change means the first arm is cheaper. H1 and H6 are \emph{refuted} (the interval excludes a saving of $15\%$); H7a, H7b, H8, H9 and H10 are \emph{supported}. The outcomes of H1 and H6 are null results by pre-registered definition, and are as reportable as the positive ones. H7b has margin $0\%$ because it is a control for a batch-size explanation, not a second saving target.
